\subsubsection{Wisconsin Old Fashioned}
\noindent
\begin{minipage}{.49\textwidth}
    \noindent\textbf{\large Wisconsin Old Fashioned\index{Wisconsin Old Fashioned}}
\end{minipage}
\hspace{.02\textwidth}
\begin{minipage}{.49\textwidth}
    \raggedleft
    \textbf{\large Brandy}\\
    \sweettartguage{-0.3}\index[drinks]{Sweet!30\%!Wisconsin Old Fashioned}
\end{minipage}
\vspace{.5em}\hrule\vspace{.5em}
\noindent
\textbf{\large Flavor Profile:} Fruity\index[flavor-profile]{Fruity}, Sweet\index[flavor-profile]{Sweet}, Raisin\index[flavor-profile]{Raisin}, Vanilla\index[flavor-profile]{Vanilla}.
\vspace{.5em}\hrule\vspace{.5em}
\noindent
\textbf{\large History:} The Wisconsin Old Fashioned was introduced in 1893 as part of the Chicago World's Fair. Brandy was getting difficult to obtain so an American liquor company called Korbel decided to make brandy and feature it at the fair and loved it so much they took it home. Not long after that the Wisconsin Old Fashioned was born. Wisconsin is the largest US consumer of brandy and makes up over 50\% of all Korbel sales.
\vspace{.5em}\hrule\vspace{.5em}
\noindent
\begin{minipage}[t]{.49\textwidth}
    \noindent
    \textbf{\large Ingredients}
    \begin{itemize}
        \item Brandy 60ml\index[ingredient]{Brandy}
        \item Orange Slices 2\index[ingredient]{Orange Slice}
        \item Maraschino Cherry 1\index[ingredient]{Maraschino Cherry}
        \item Sugar Cube 1\index[ingredient]{Sugar Cube}
        \item Angostura Bitters 3 dashes\index[ingredient]{Angostura Bitters}
    \end{itemize}
    \noindent
    \textbf{\large Top off}
    \begin{itemize}
        \item 7-up\index[ingredient]{7-Up} or Squirt\index[ingredient]{Squirt}
    \end{itemize}
    \noindent
    \textbf{\large Garnish}
    \begin{itemize}
        \item Cherry\index[ingredient]{Cherry}
        \item Orange Slice\index[ingredient]{Orange Slice}
    \end{itemize}
\end{minipage}
\hspace{.01\textwidth}\vrule\hspace{.01\textwidth}
\begin{minipage}[t]{.49\textwidth}
    \noindent
    \textbf{\large Tools/Glassware}
    \begin{itemize}
        \item Jigger
        \item Bar Spoon
        \item Muddle
        \item Double Rocks Glass (360ml)\index[glass]{Double Rocks Glass}
    \end{itemize}
\end{minipage}
\vspace{.5em}\hrule\vspace{.5em}
\noindent
\begin{minipage}{\textwidth}
    \noindent
    \textbf{\large Instructions}
    \begin{itemize}
        \item Place orange slices and cherry into glass.
        \item Add sugar cube and 3 dashes of Angostura bitters.
        \item Lightly muddle everything until sugar is dissolved.
        \item Add ice to glass, pour in Brandy and stir.
        \item Add 7-up or Squirt to taste.
        \item Garnish with orange slice and cherry.
    \end{itemize}
\end{minipage}
\vspace{.5em}\hrule\vspace{.5em}
\noindent
\textbf{\large Helpful Tip:} Add 7-up for "sweet", Squirt for "sour", or soda water for "press" style old fashioned.
\newpage
